% ============================================================
% S2 导学件量产·批3 件2 —— 1.2.1 前几何衔接位（衔接1～29 独立节页）
% 轮4 成卷·S2 批3｜2026-09-12｜生成链：xelatex main.tex（两遍）→ python render.py（150dpi → png/pageN.png）
%
% 内容源＝`成卷/题面库/衔接节.md`（29 题已冻＋五组排布序＋组内序约束＋四心填空讲块[190]–[214]＋点拨段[301]）；
%   排布硬约束＝全章定稿汇总 §3.2/§3.3（组内序四条：①衔接18 先于衔接24/29（张角结论前置，源号 -3→-9/-14）；
%   ②衔接20 与衔接22(2) 同页推导（等面积/切线长「公式—数值」讲一次，源号 -5 与 -7(2) 同式约束对组二侧
%   即衔接5 与衔接7(2) 同页）；③垂心族(衔接17/23/25)与两心重合族(衔接21/28)分节排布勿邻接堆叠；
%   ④图指代用题必配图（衔接2/3 等，题面【图】21 处全带）——成品版面按 衔接1→29 升序即满足组序，
%   编译后跑序断言（R7 防「白做29题」）见批3值台账。
% 件样式＝承练习件版式（练习线 分区组行\zuhang＋悬挂题号＋选项件 \lxopt 局部宏照抄 v11 练习线/main.tex；
%   四心讲块＝全品课前预习式扩展：\zsd 同构标题宏局部件＋\tiaomu 条目）；六宏＝原样拷贝零改动。
% 图＝源图位图（禁绘三维）：figs/ 28 件自 衔接件（29题）.docx word/media/ 原件解包未改绘
%   （题面 21 处＋讲块 7 处；源 extent 微小（5–9mm 档）＝源件版式病，成品按像素比例放大至可读宽，登记）。
%   【目验修红0913】装配本逐页目验裁：衔接5 题首「如图，」删语（题面文字自足可解，先例＝拓-046/047/058）。
%     惟本项与其余四处不同：image15 实有配图且渲染正常（成品 p17 右栏首），缺陷＝图题异栏相距过远非「无图」；
%     删语后 image15 照印作示意插图、无悬空指代，图 21 处计数不变，「如图」指代降为 20 处（§3.2④ 必配图口径不受影响）。
% 页码口径：本件为导学件册内 1.2.1 前插页，页脚件名＝导学件；跨本连续页码＝装订轮。
% 答案同号段：答案册同节同号 衔接-1～衔接-29（排产单 §2.1/§五.5）；版面题号印式「衔接N.」＝号段映射登记台账。
% 编译门：三 0——读数见批3值台账。
% ============================================================
\documentclass[fontset=none]{ctexart}
\usepackage{qp-m3} % 换装：六宏→qp-m3 单包（钉后版 md5 7c3930362be8a0a2bdf21bbf8ac16573，两补钉随版）

% ================= 页脚参数宏 =================
\renewcommand{\qpzhangming}{第一章\quad 空间向量与立体几何}
\renewcommand{\qpjianming}{导学件}
\renewcommand{\qpceming}{高中数学\quad 选择性必修第一册(人教B版)} % 换装回钉：qp-m3 默认册名＝选必二，防页脚漂移（试迁规约）

% ================= 局部宏（承练习线冻结件形制，公共模块语义不动） =================
% JT-06 分区组行（照抄 练习线/main.tex 定义）——本件作五组组行
\newcommand{\zuhang}[2]{\par\glueguard{2}\addvspace{1.5mm}%
  {\noindent{\fontsize{10.5pt}{18pt}\selectfont\heibf #1}{\fontsize{10.5pt}{18pt}\selectfont：#2}\par}%
  \addvspace{1.3mm}}
% TJ-07 选项段（照抄 练习线/main.tex 定义；行距＝例题区档 19.2pt）
\renewcommand{\lxopt}[1]{\par\noindent\hangindent=\xhang\hangafter=0
  {\fontsize{10.5pt}{\qpTJgLeadLiti}\selectfont #1\par}}
% 衔接悬挂题号（TJ-02 同构：题号「衔接N.」黑体加粗坐栏左缘，单号 11.0mm／双位 13.0mm）
\newlength{\xhang}
\newcommand{\xthao}[1]{\par\glueguard{1}%
  \ifnum#1>9\setlength{\xhang}{13.0mm}\else\setlength{\xhang}{11.0mm}\fi
  \noindent\hangindent=\xhang\hangafter=0
  \llap{\hbox to\xhang{{\heibf 衔接}\textbf{#1.}\hss}}}
% 小题行（同悬挂续行）
\newcommand{\subl}[1]{\par\noindent\hangindent=\xhang\hangafter=0 #1\par}
% 源图位图居中定行（承导学件 \bindp 图式）
\newcommand{\fig}[2][34.0mm]{\par\vspace{1.9mm}\penalty10000\noindent\makebox[\linewidth][c]{\includegraphics[width=#1]{figs/#2.png}}\par\vspace{-1.0mm}\penalty10000}
% 知识讲块标题（\zsd 同构件：◆＋12.03pt \heibf，仅去「知识点N」编号——讲块不占知识点号段）
\newcommand{\zhublock}[1]{\par\addvspace{6.3pt}\glueguard{4}%
  \noindent{\fontsize{12.03pt}{17pt}\selectfont\heibf \huafu\hspace{2.4mm}#1}\par\addvspace{2pt}}
% 点拨行（源「点拨 　…」仿宋件）
\newcommand{\diangu}[1]{\par\glueguard{2}\addvspace{1.2mm}\noindent{\fangsong 【点拨】#1}\par\addvspace{1.2mm}}

\begin{document}
%成书重装五本跨本连续页码（G7方案B，承M2/M3装配先例）setcounter 注入位
\setcounter{page}{1}

% ============================================================
% 衔接节页（1.2.1 前几何衔接位；五组排布＋四心讲块置第四组之前＋点拨段[301]原样＋方法链接三挂点）
% ============================================================
\keshi{衔接\quad 初等几何必会（1.2.1 前）}

\vspace{1.0mm}
{\fangsong 本节复习 1.2.1《空间中的点、直线与空间向量》必用的初等几何知识，共 29 题，独立编号衔接1—衔接29（答案册同节同号）．\par}

\vspace{1.2mm}
\begin{multicols}{2}
\emergencystretch=1em

% ---------------- 第一组：平行线分线段成比例／相似／中位线（衔接1～5） ----------------
\zuhang{一、平行线分线段成比例、相似与中位线}{本组 5 题(衔接必会)}

\xthao{1}\tieside{衔接必会}如图，在\(\triangle ABC\)中，\(\frac{AD}{AB}=\frac{1}{4}\)，\(\frac{AE}{AC}=\frac{1}{2}\)，\(BE\)和\(CD\)相交于点\(F\)，则\(\frac{DF}{DC}=\)\kongbai{}.
\fig[30.0mm]{image8}

\begin{ansblock}[导-衔接-1]
% ans:导-衔接-1
\ansitem{1}{\(\dfrac{3}{7}\)}
\end{ansblock}

\xthao{2}\tieside{衔接必会}如图，已知\(M\)为平行四边形\(ABCD\)的边\(AB\)的中点，\(CM\)交\(BD\)于点\(E\)，则图中阴影部分的面积与平行四边形\(ABCD\)面积的比是\kongbai{}.
\fig[26.0mm]{image10}

\begin{ansblock}[导-衔接-2]
% ans:导-衔接-2
\ansitem{2}{\(1:3\)}
\end{ansblock}

\xthao{3}\tieside{衔接必会}如图，\(D\)，\(E\)分别为边\(AB\)，\(AC\)的中点，则\(\frac{BP}{BE}\)的值\kongbai{}.
\fig[45.0mm]{image11}

\begin{ansblock}[导-衔接-3]
% ans:导-衔接-3
\ansitem{3}{\(\dfrac{2}{3}\)}
\end{ansblock}

\xthao{4}\tieside{衔接必会}如图，\(BD\)、\(CE\)是\(\triangle ABC\)的中线，\(P\)、\(Q\)分别是\(BD\)、\(CE\)的中点，则\(PQ:BC\)等于\kongbai{}.
\fig[34.0mm]{image13}

\begin{ansblock}[导-衔接-4]
% ans:导-衔接-4
\ansitem{4}{\(1:4\)}
\end{ansblock}

\xthao{5}\tieside{衔接必会}四边形\(ABCD\)的对角线相交于点\(O\)，\(\angle BAC=\angle CDB\)，求证：\(\angle DAC=\angle CBD\).
\fig[34.0mm]{image15}

% ---------------- 第二组：射影定理组（衔接6～11） ----------------
\begin{ansblock}[导-衔接-5]
% ans:导-衔接-5
\ansitem{5}{证明见解析}
\ansnote{解析}{\(\angle OAB=\angle ODC\)＋对顶角 \(\Rightarrow\triangle OAB\sim\triangle ODC\) \(\Rightarrow OA:OD=OB:OC\Rightarrow\triangle OAD\sim\triangle OBC\Rightarrow\angle DAC=\angle CBD\)。}
\end{ansblock}

\zuhang{二、射影定理组}{本组 6 题(衔接必会)}

\xthao{6}\tieside{衔接必会}在\(\mathrm{Rt}\triangle ABC\)中，\(\angle BAC=90^{\circ}\)，\(AD\perp BC\)于点\(D\)，若\(\frac{AC}{AB}=\frac{3}{4}\)，则\(\frac{AD}{AC}=\)\kongbai{}，\(\frac{BD}{CD}=\)\kongbai{}.
\fig[38.0mm]{image16}

\begin{ansblock}[导-衔接-6]
% ans:导-衔接-6
\ansitem{6}{\(\dfrac{4}{5}\)；\(\dfrac{16}{9}\)}
\end{ansblock}

\xthao{7}\tieside{衔接必会}如图，在\(\triangle ABC\)中，\(\angle ACB=90^{\circ}\)，\(CD\perp AB\)，\(AD=6\)，\(CD=2\sqrt{3}\)，求：
\fig[30.0mm]{image17}
\subl {(1)\(\angle A\)的度数；\quad (2)\(\triangle ABC\)的面积.}

\begin{ansblock}[导-衔接-7]
% ans:导-衔接-7
\ansitem{7}{(1)\(30^{\circ}\)；(2)\(8\sqrt{3}\)}
\end{ansblock}

\xthao{8}\tieside{衔接必会}一个直角三角形的一条直角边为\(3\)，斜边上的高为\(2.4\)，则这个直角三角形的面积为(\kongwei)
\lxopt{\makebox[\dimexpr0.5\linewidth-0.5\xhang-1em\relax][l]{A．\(7.2\)；}\makebox[\dimexpr0.5\linewidth-0.5\xhang-1em\relax][l]{B．\(6\)；}}
\lxopt{\makebox[\dimexpr0.5\linewidth-0.5\xhang-1em\relax][l]{C．\(12\)；}\makebox[\dimexpr0.5\linewidth-0.5\xhang-1em\relax][l]{D．\(24\)}}

\begin{ansblock}[导-衔接-8]
% ans:导-衔接-8
\ansitem{8}{B（\(6\)）}
\end{ansblock}

\xthao{9}\tieside{衔接必会}已知在梯形\(ABCD\)中，\(DC\parallel AB\)，\(\angle D=90^{\circ}\)，\(AC\perp BC\)，\(AB=10\)，\(AC=6\)，则此梯形的面积为\kongbai{}.
\fig[40.0mm]{image19}

\begin{ansblock}[导-衔接-9]
% ans:导-衔接-9
\ansitem{9}{\(816/25\)}
\end{ansblock}

\xthao{10}\tieside{衔接必会}已知直角三角形周长为\(48\,\mathrm{cm}\)，一锐角平分线分对边为\(3:5\)两部分.
\fig[24.0mm]{image20}
\subl {(1)求直角三角形的三边长；(2)求两直角边在斜边上的射影的长.}

\begin{ansblock}[导-衔接-10]
% ans:导-衔接-10
\ansitem{10}{(1)\(20\)，\(12\)，\(16\) cm；(2)\(\dfrac{36}{5}\)，\(\dfrac{64}{5}\) cm}
\end{ansblock}

\xthao{11}\tieside{衔接必会}如图，在\(\triangle ABC\)中，\(AD\)为\(BC\)边上的高，过\(D\)作\(DE\perp AB\)，\(DF\perp AC\)，\(E\)，\(F\)为垂足.求证：
\fig[38.0mm]{image22}
\subl {(1)\(AE\cdot AB=AF\cdot AC\)；(2)\(\triangle AEF\backsim\triangle ACB\).}

% ---------------- 第三组：内角／外角平分线定理（衔接12～15） ----------------
\begin{ansblock}[导-衔接-11]
% ans:导-衔接-11
\ansitem{11}{证明见解析}
\ansnote{解析}{双射影：\(AE\cdot AB=AD^{2}=AF\cdot AC\)；(2) 共角夹边成比例逆相似。}
\end{ansblock}

\zuhang{三、三角形的内角与外角平分线定理}{本组 4 题(衔接必会)}

\xthao{12}\tieside{衔接必会}已知，如图所示，\(AD\)是\(\triangle ABC\)的内角\(\angle BAC\)的平分线，求证：\(\frac{BA}{AC}=\frac{BD}{DC}\).
\fig[34.0mm]{image6}

\begin{ansblock}[导-衔接-12]
% ans:导-衔接-12
\ansitem{12}{证明见解析}
\ansnote{解析}{内角平分线定理：过 \(C\) 作平行线法（面积法同源）。}
\end{ansblock}

\xthao{13}\tieside{衔接必会}如图，在\(\triangle ABC\)中，\(AD\)是角\(BAC\)的平分线，\(AB=5\,\mathrm{cm}\)，\(AC=4\,\mathrm{cm}\)，\(BC=7\,\mathrm{cm}\)，求\(BD\)的长.
\fig[32.0mm]{image25}

\begin{ansblock}[导-衔接-13]
% ans:导-衔接-13
\ansitem{13}{\(\dfrac{35}{9}\)}
\end{ansblock}

\xthao{14}\tieside{衔接必会}已知，如图所示，\(AD\)是\(\triangle ABC\)中\(\angle BAC\)的外角\(\angle CAF\)的平分线，求证：\(\frac{BA}{AC}=\frac{BD}{DC}\).
\fig[38.0mm]{image7}

\begin{ansblock}[导-衔接-14]
% ans:导-衔接-14
\ansitem{14}{证明见解析}
\ansnote{解析}{外角平分线定理；合比定理：\(BD:DC=AB:AC\Rightarrow BD:(BD+DC)=AB:(AB+AC)\)。}
\end{ansblock}

\xthao{15}\tieside{衔接必会}在\(\triangle ABC\)中，\(AD\)是\(\angle BAC\)的平分线，\(AB-AC=5\)，\(BD-CD=3\)，\(DC=8\)，则\(AB=\)\kongbai{}.
\fig[40.0mm]{image27}

% ---------------- 四心填空讲块（源元素[190]–[214]，置第四组之前；全品课前预习式扩展） ----------------
\begin{ansblock}[导-衔接-15]
% ans:导-衔接-15
\ansitem{15}{\(\dfrac{55}{3}\)}
\end{ansblock}

\zhublock{知识讲块\quad 三角形的“四心”}

\tiaomuz{1}{三角形的重心：(1)三角形的三条中线相交于一点，这个交点称为三角形的重心.\zhuzhu{点\(D\)、\(E\)、\(F\)是中点，则点\(O\)是\(\triangle ABC\)的重心.}}
\fig[36.0mm]{image28}
\bindp (2)三角形的重心在三角形的内部，恰好是每条中线的三等分点．如上图，\(AO=2OD\)，\(CO=2OE\)，\(BO=2OF\)．\par\addvspace{0.80mm}
\tiaomutail

\tiaomuz{2}{三角形的内心：(1)三角形的三条角平分线相交于一点，是三角形的内心.\zhuzhu{\(AD\)、\(CE\)、\(BF\)是角平分线，则点\(O\)是\(\triangle ABC\)的内心.}}
\fig[34.0mm]{image29}
\bindp (2)三角形的内心在三角形的内部，它到三角形的三边的距离相等.
\fig[38.0mm]{image30}\tiaomutail

\tiaomuz{3}{三角形的垂心：(1)三角形的三条高所在直线相交于一点，该点称为三角形的垂心.}
\fig[38.0mm]{image31}
\bindp (2)锐角三角形的垂心一定在三角形的内部，直角三角形的垂心为他的直角顶点，钝角三角形的垂心在三角形的外部.
\fig[60.0mm]{image32}\tiaomutail

\tiaomuz{4}{三角形的外心：(1)三角形三边的垂直平分线所在直线相交于一点，该点称为三角形的外心.\zhuzhu{\(OE\)、\(OD\)、\(OF\)分别是\(AB\)、\(BC\)、\(AC\)的垂直平分线，点\(O\)是三角形的外心.}}
\fig[40.0mm]{image33}
\bindp (2)三角形的外心到三个顶点的距离相等.
\fig[38.0mm]{image34}\tiaomutail

\tiaomu{5}{拓展：(1)等腰三角形底边上三线（角平分线、中线、高线）合一.因而在等腰三角形\(ABC\)中，三角形的内心\(I\)、重心\(G\)、垂心\(H\)必然在一条直线上.(2)正三角形三条边长相等，三个角相等，且四心（内心、重心、垂心、外心）合一，该点称为正三角形的中心.}

% ---------------- 第四组：三角形四心性质证明（衔接16～20） ----------------
\zuhang{四、三角形“四心”的性质证明}{本组 5 题(衔接必会)}

\xthao{16}\tieside{衔接必会}求证\(D\)、\(E\)、\(F\)分别为\(\triangle ABC\)三边\(BC\)、\(CA\)、\(AB\)的中点，\(AD\)、\(BE\)、\(CF\)交于一点，且三条中线都被该点分成\(2:1\).
\fig[36.0mm]{image35}

\begin{ansblock}[导-衔接-16]
% ans:导-衔接-16
\ansitem{16}{证明见解析}
\ansnote{解析}{重心三等分（向量/中位线双链）；取 \(AC\) 中点。}
\end{ansblock}

\xthao{17}\tieside{衔接必会}已知\(\triangle ABC\)中，\(AD\perp BC\)于\(D\)，\(BE\perp AC\)于\(E\)，\(AD\)与\(BE\)交于\(H\)点，求证\(CH\perp AB\).
\fig[34.0mm]{image36}

\begin{ansblock}[导-衔接-17]
% ans:导-衔接-17
\ansitem{17}{证明见解析}
\ansnote{解析}{双高 \(\Rightarrow H\)、\(D\)、\(C\)、\(E\) 与 \(A\)、\(B\)、\(D\)、\(E\) 两组四点共圆，转角即得 \(CH\) 延长线\(\perp AB\)。}
\end{ansblock}

\xthao{18}\tieside{衔接必会}已知：\(\triangle ABC\)是三边都不相等的三角形，点\(O\)和点\(P\)是这个三角形内部两点.
\subl {(1)如图①，如果点\(P\)是这个三角形三个内角平分线的交点，那么\(\angle BPC\)和\(\angle BAC\)有怎样的数量关系？请说明理由；}
\subl {(2)如图②，如果点\(O\)是这个三角形三边垂直平分线的交点，那么\(\angle BOC\)和\(\angle BAC\)有怎样的数量关系？请说明理由；}
\subl {(3)如图③，如果点\(P\)（三角形三个内角平分线的交点），点\(O\)（三角形三边垂直平分线的交点）同时在不等边\(\triangle ABC\)的内部，那么\(\angle BPC\)和\(\angle BOC\)有怎样的数量关系？请直接回答.}
\fig[84.0mm]{image37}

\begin{ansblock}[导-衔接-18]
% ans:导-衔接-18
\ansitem{18}{(1)\(\angle BPC=90^{\circ}+\dfrac{1}{2}\angle BAC\)；(2)\(\angle BOC=2\angle BAC\)；(3)\(4\angle BPC-\angle BOC=360^{\circ}\)}
\end{ansblock}

\xthao{19}\tieside{衔接必会}已知\(\triangle ABC\)的三边长分别为\(BC=a\)，\(AC=b\)，\(AB=c\)，\(I\)为\(\triangle ABC\)的内心，且\(I\)在\(\triangle ABC\)的边\(BC\)，\(AC\)，\(AB\)上的射影分别为\(D\)，\(E\)，\(F\)，求证：\(AE=AF=\frac{b+c-a}{2}\).
\fig[40.0mm]{image39}

\begin{ansblock}[导-衔接-19]
% ans:导-衔接-19
\ansitem{19}{证明见解析}
\ansnote{解析}{切线长相等三元一次 \(\Rightarrow AE=AF=\dfrac{b+c-a}{2}\)。}
\end{ansblock}

\xthao{20}\tieside{衔接必会}若三角形\(ABC\)的面积为\(S\)，且三边长分别为\(a\)、\(b\)、\(c\)，求三角形的内切圆的半径.

% ---------------- 第五组：三角形四心运用（衔接21～29） ----------------
\begin{ansblock}[导-衔接-20]
% ans:导-衔接-20
\ansitem{20}{\(r=\dfrac{2S}{a+b+c}\)}
\end{ansblock}

\zuhang{五、三角形“四心”的运用}{本组 9 题(衔接必会)}

\xthao{21}\tieside{衔接必会}若已知\(O\)为三角形\(ABC\)的重心和内心，求证三角形\(ABC\)为等边三角形.

\begin{ansblock}[导-衔接-21]
% ans:导-衔接-21
\ansitem{21}{证明见解析}
\ansnote{解析}{重心＝内心 \(\Rightarrow\) 中线即角平分线 \(\Rightarrow\) 逐向等腰 \(\Rightarrow\) 等边。}
\end{ansblock}

\xthao{22}\tieside{衔接必会}在\(\triangle ABC\)中，\(AB=AC=3\)，\(BC=2\)，求
\subl {(1)\(\triangle ABC\)的面积\(S_{\triangle ABC}\)及\(AC\)边上的高\(BE\)；(2)\(\triangle ABC\)的内切圆的半径\(r\)；(3)\(\triangle ABC\)的外接圆的半径\(R\).}

\begin{ansblock}[导-衔接-22]
% ans:导-衔接-22
\ansitem{22}{(1)\(S=2\sqrt{2}\)，\(BE=\dfrac{4\sqrt{2}}{3}\)；(2)\(r=\dfrac{\sqrt{2}}{2}\)；(3)\(R=\dfrac{9\sqrt{2}}{8}\)}
\ansnote{解析}{(2) 即用衔接-20 公式；(3) \(R=\dfrac{abc}{4S}\)。}
\end{ansblock}

\xthao{23}\tieside{衔接必会}在\(\triangle ABC\)中，\(AB=10\)，\(AD\)是\(BC\)边上的高，已知\(AD=BD\)，点\(H\)、\(O\)分别是\(\triangle ABC\)的垂心和外心，则\(HO\)的最小长度为\kongbai{}.
\fig[32.0mm]{image46}
\diangu{建立直角坐标系，把几何问题转化为代数问题（把\(OH\)线段的最小值化为代数式\(\frac{5}{2}(m-2\sqrt{2})^{2}+5\)的最小值），这是解析几何的方法，在高中经常用到.}

\begin{ansblock}[导-衔接-23]
% ans:导-衔接-23
\ansitem{23}{\(\sqrt{5}\)}
\end{ansblock}

\xthao{24}\tieside{衔接必会}\(O\)是\(\triangle ABC\)的内切圆圆心，\(\angle C=90^{\circ}\)，\(\angle BOC=105^{\circ}\)，\(BC=20\,\mathrm{cm}\)，则\(AC\)长为\kongbai{}\(\mathrm{cm}\).

\begin{ansblock}[导-衔接-24]
% ans:导-衔接-24
\ansitem{24}{\(20\sqrt{3}\)}
\end{ansblock}

\xthao{25}\tieside{衔接必会}在\(\triangle ABC\)中，\(\angle A\)是钝角，\(H\)是垂心，\(AH=BC\)，则\(\angle BHC=\)\kongbai{}.

\begin{ansblock}[导-衔接-25]
% ans:导-衔接-25
\ansitem{25}{\(45^{\circ}\)}
\ansnote{解析}{\(AH=2R|\cos A|=BC=2R\sin A\)（\(A\) 钝）\(\Rightarrow A=135^{\circ}\Rightarrow\angle BHC=180^{\circ}-A\)。}
\end{ansblock}

\xthao{26}\tieside{衔接必会}在\(\triangle ABC\)中，两中线\(AD\)与\(CF\)相交于点\(G\)，若\(\angle AFC=45^{\circ}\)，\(\angle AGC=60^{\circ}\)，则\(\angle ACF\)的度数为\kongbai{}.

\begin{ansblock}[导-衔接-26]
% ans:导-衔接-26
\ansitem{26}{\(75^{\circ}\)}
\end{ansblock}

\xthao{27}\tieside{衔接必会}如图，在\(\mathrm{Rt}\triangle ABC\)中，\(\angle BCA=90^{\circ}\)，中线\(CM\)垂直于中线\(BN\)，且\(BC=2\)，求\(BN\)的长.
\fig[42.0mm]{image51}

\begin{ansblock}[导-衔接-27]
% ans:导-衔接-27
\ansitem{27}{\(\sqrt{6}\)}
\end{ansblock}

\xthao{28}\tieside{衔接必会}求证：若三角形的垂心和重心重合，求证：该三角形为正三角形．

\begin{ansblock}[导-衔接-28]
% ans:导-衔接-28
\ansitem{28}{证明见解析}
\ansnote{解析}{垂心＝重心 \(\Rightarrow\) 中线\(\perp\)对边 \(\Rightarrow\) 等腰循环 \(\Rightarrow\) 正三角形。}
\end{ansblock}

\xthao{29}\tieside{衔接必会}锐角\(\triangle ABC\)中，\(AB>BC>CA\)，\(O\)、\(I\)、\(H\)分别是它的外心、内心、垂心，已知\(\angle A=60^{\circ}\)，求证：
\subl {(1)\(\angle OIH-\angle ABC\)是一个定值；(2)\(\angle OIH+\angle ACB\)也是一个定值.}
\fig[40.0mm]{image54}

% ---------------- 方法链接（汇总 §3.3 三挂点逐一点名；点拨段[301]原样已随衔接23） ----------------
\begin{ansblock}[导-衔接-29]
% ans:导-衔接-29
\ansitem{29}{证明见解析}
\ansnote{解析}{五点共圆链；定值 (1)\(120^{\circ}\) (2)\(240^{\circ}\)（\(A=60^{\circ}\) 条件下）。}
\end{ansblock}

\zhublock{方法链接\quad 本章的三处接口}

\bindp ①等面积法：衔接20 的“\(r=\frac{2S}{a+b+c}\)”（衔接22(2) 即用此式算数）到本章升维为内切球“\(r=\frac{3V}{S_{\text{全}}}\)”——1.2.5 的“体积＋内切球”题组即该法的空间用场；

\bindp ②张角与外接圆半径：衔接18(2) 的外心张角“\(\angle BOC=2\angle A\)”与衔接22(3) 求\(R\)（勾股／方程，\(R^{2}=\)弦心距\({}^{2}+\)半弦\({}^{2}\)）直通本章核心算法“外接球球心在过底面外心且垂直底面的轴上”——平面圆→空间球直接升维；

\bindp ③坐标法：衔接23 以\(D\)为原点、\(BC\) 所在直线与高线\(DA\)为两轴建系求\(OH\)最值，正是本章 1.1.3 空间直角坐标系的预告——衔接23 后点拨“解析几何的方法，在高中经常用到”即此挂点.\par\addvspace{0.80mm}

% ============================================================
% 衔接节升格·结构挂位点（换装正装轮0914 立·义务总表§一.1②）
%   全制式补＝考点探究（\huaxing{课}{中}{探}{究}{考点探究\quad 素养小结}＋\tjdnr/\liB 例变骨架）
%           ＋课堂评价 G5（\begin{ketang}{本组共5题，1—3为单选，4—5为填空．}…\end{ketang}，
%             判/选/填 恒5 内构）——补题命制挂后续轮；本轮先立结构不印空节（防空节版面红旗），
%             qp-m3 制式宏已就位，命制到件时按课时件批1—6 样板展开即入位。
% ============================================================
\end{multicols}

\tailfill
\end{document}
